% KKchemstruct.sty を同じディレクトリか TeX の検索パスへ配置し、
% LuaLaTeX で2回コンパイルしてください。
\documentclass[a4paper,11pt]{ltjsarticle}
\usepackage[margin=22mm]{geometry}
\usepackage{KKchemstruct}
\usepackage[most]{tcolorbox}
\pagestyle{empty}

\definecolor{qaccent}{HTML}{0EA5E9}      % sky-500
\definecolor{qaccentDark}{HTML}{0369A1}  % sky-700
\definecolor{qsoft}{HTML}{F0F9FF}        % sky-50
\definecolor{qtext}{HTML}{1E293B}        % slate-800

\newtcolorbox{ExamplePanel}[1]{
  enhanced,
  colback=qsoft,
  colframe=qaccent,
  boxrule=.6pt,
  arc=1.5mm,
  left=5mm,
  right=5mm,
  top=4mm,
  bottom=4mm,
  title={#1},
  fonttitle=\sffamily\bfseries,
  coltitle=white,
}

\newcommand{\pagetitle}[1]{%
  \begin{center}
    {\sffamily\bfseries\Large\color{qaccentDark}#1}\par
    \vspace{1mm}
    {\color{qaccent}\rule{.72\linewidth}{.7pt}}
  \end{center}
  \vspace{5mm}}

\begin{document}

% Figure 1: quick start
\pagetitle{KKchemstruct クイックスタート}

\begin{ExamplePanel}{置換ベンゼン}
  \begin{center}
    \KKname{\KKsalicylicacid}{サリチル酸}
    \hspace{5em}
    \KKname{\KKaspirin}{アセチルサリチル酸}
  \end{center}
\end{ExamplePanel}

\vspace{5mm}

\begin{ExamplePanel}{示性式とオキソ酸}
  \begin{center}
    \KKname{\KKglycerol}{グリセリン}
    \hspace{4em}
    \KKname{\KKsulfuricacid}{硫酸}
    \hspace{4em}
    \KKname{\KKphosphoricacid}{リン酸}
  \end{center}
\end{ExamplePanel}

\vspace{5mm}
\begin{center}
  \small\ttfamily
  \string\KKsalicylicacid\qquad
  \string\KKglycerol\qquad
  \string\KKsulfuricacid
\end{center}

\newpage

% Figure 2: benzene and chain
\pagetitle{ベンゼン環と示性式}

\begin{ExamplePanel}{置換位置は真上を 1 として時計回り}
  \begin{center}
    \KKname{\KKbenzene[2=OH,3=COOH]}{2=OH, 3=COOH}
    \hspace{4em}
    \KKname{\KKbenzene[1=COOH,4=COOH]}{1=COOH, 4=COOH}
    \hspace{4em}
    \KKname{\KKbenzene[]}{無置換}
  \end{center}
\end{ExamplePanel}

\vspace{5mm}

\begin{ExamplePanel}{環の調整と横向きの環}
  \begin{center}
    \KKname{\KKbenzene[1=OH,kekule=b]}{kekule=b}
    \hspace{5em}
    \KKname{\KKbenzene[2=OH,angle=30,sep=1.0em]}{angle=30, sep=1.0em}
  \end{center}

  \medskip
  本文中にも \KKphenylring{} や \KKphenyl{COOH} を自然に置けます。
\end{ExamplePanel}

\vspace{5mm}

\begin{ExamplePanel}{示性式は見える結合線の長さをそろえる}
  \begin{center}
    \KKname{\KKchemchain{R_1-C\KKcarbonyl-OH}}{カルボン酸の一般形}
    \hspace{5em}
    \KKname{\KKtriglyceride{R_1}{R_2}{R_3}}{油脂}
  \end{center}
\end{ExamplePanel}

\newpage

% Figure 3: reactions
\pagetitle{反応式}

\begin{ExamplePanel}{\texttt{\string\KKscheme}: 1 本の反応矢印}
  \KKscheme{%
    \KKname{\KKsalicylicacid}{サリチル酸}
    \arrow{->[\ck{CH_3OH}][エステル化]}
    \KKname{\KKmethylsalicylate}{サリチル酸メチル}}
\end{ExamplePanel}

\vspace{8mm}

\begin{ExamplePanel}{\texttt{\string\KKbranchscheme}: 2 本へ枝分かれ}
  \KKchemsetup{compound sep=3.8em,branch coeff=2.2}
  \KKbranchscheme
    {\KKname{\KKsalicylicacid}{サリチル酸}}
    {\ck{CH_3OH}}{エステル化}
    {\KKname{\KKmethylsalicylate}{サリチル酸メチル}}
    {\ck{(CH_3CO)_2O}}{アセチル化}
    {\KKname{\KKaspirin}{アセチルサリチル酸}}
\end{ExamplePanel}

\vspace{5mm}
\begin{center}
  \small
  矢印ラベルの化学式は \texttt{\string\ck\{CH\_3OH\}} のように書きます。
\end{center}

\newpage

% Figure 4: annotations and sizing
\pagetitle{注釈と寸法変更}

\begin{ExamplePanel}{注釈}
  \begin{center}
    \KKchemchain{R-C\KKcarbonyl-@{leavingA}OH}\ +\
    \KKchemchain{@{leavingB}H-O-R'}%
    \KKchemmark{leavingA}{leavingB}

    \vspace{9mm}

    \KKbenzene[2=@{hydrogenO}OH,3=C@{carbonylO}OOH]%
    \KKhbond{hydrogenO}{carbonylO}
    \hspace{6em}
    \KKchemcross{\KKphenol}
  \end{center}
\end{ExamplePanel}

\vspace{7mm}

\begin{ExamplePanel}{\texttt{\string\KKchemsetup} はグループ内だけにも適用できる}
  \begin{center}
    \KKname{\KKsalicylicacid}{既定値}
    \hspace{6em}
    \begingroup
      \KKchemsetup{ring sep=2.0em,bond width=0.09em,double bond sep=0.4em}
      \KKname{\KKsalicylicacid}{変更後}
    \endgroup
  \end{center}
\end{ExamplePanel}

\vspace{5mm}
\begin{center}
  \small
  \texttt{\string\KKchemmark} と \texttt{\string\KKhbond} を使う文書は 2 回コンパイルします。
\end{center}

\end{document}
